\section{XAI Methods}\label{sec:methods_xai}

In this section, we leverage the neural network decomposition framework introduced in the previous sections to boost two common XAI attribution methods: LayerCAM and LRP.

\subsection{SplitCAM}\label{subsec:split_cam}

Class Activation Mapping (CAM) \citep{Zhou2016} and its variants (Grad-CAM \citep{Selvaraju2017}, Grad-CAM++ \citep{Chattopadhay2018}, LayerCAM \citep{Jiang2021LayerCAM}) generate visual explanations by combining layer activation maps with class specific weights or layer gradients.

For layer $l$, LayerCAM computes the saliency map as
\begin{equation} \label{eq:layercam}
    \text{LayerCAM}^{(l)} = \text{ReLU}\left(\sum_{c} \frac{\partial y}{\partial a^{(l)}_c} \odot a^{(l)}_c\right)
\end{equation}
where $a^{(l)}_c$ is the activation map of channel $c$ at layer $l$, $y$ is the class score, and $\odot$ denotes element-wise multiplication.

\emph{SplitCAM} extends the LayerCAM framework to our split-pair architecture.
For a given layer~$l$ in the decomposed model we combine the two split-stream activations with its gradients, namely $\delta_{\text{shift}}^{(l,+,g)}$ (one of the shifted local sensitivities) and $\delta_{\text{shift}}^{(l,g)}$ (a combination of such) defined in \cref{sec:backward}. Different from the original LayerCAM formula \cref{eq:layercam}, SplitCAM omits the positive filtering by the ReLU function.

\begin{align}
    \text{SplitCAM}^{(l, +, g)} &\coloneqq \sum_{c} \delta_{\text{shift}}^{(l,+,g)} \odot a^{(l,+)}_c, \\
    \text{SplitCAM}^{(l, g)} &\coloneqq \sum_{c} \delta_{\text{shift}}^{(l,g)} \odot a^{(l)}_c,
\end{align}
where $a^{(l,+)}_c$ denotes the positive activation map of channel~$c$.
$\text{SplitCAM}^{(l, -, g)}$ and the variants for $h$ are defined analogously.

In the forward pass, we scale the entire activation tensor by a fixed \emph{minimize factor} $\theta$ whenever the absolute value of any activation entry exceeds a predefined \emph{minimize threshold} $\Theta$.
In all our experiments we use $\theta = 0.1$ and $\Theta = 10$.
We apply corrective shifts to the split-pair activations such that their difference matches the original network's activation.

In the backward pass we apply the shifting procedure described in \cref{sec:backward} and slightly correct the four sensitivity maps in each layer to ensure the backward invariant \cref{eq:backward_invariant}.

\paragraph{SplitGrad} As a simple attribution method we also propose SplitGrad, which resembles the channel mean of the shifted split-pair gradients in the previous setup.

\subsection{SplitLRP}\label{subsec:split_lrp}

Layer-wise Relevance Propagation \citep{Bach2015} redistributes the output score iteratively from the
output layer to the input layer depending on the weights and activations. There are several extensions to the classical LRP rule. 
Here we present the $\gamma$-rule, introduced by Montavon et al.\ in \citep{Montavon_2022_gammaLRP}, that tries to emphasize positive influences to the network's prediction by increasing the influence of positive weights.

\paragraph{LRP-\(\gamma\) rule (linear / convolution layers)}
For an affine mapping \(z^{(l+1)}_i=\sum_j W^{(l+1)}_{ij} a^{(l)}_j + b^{(l+1)}_i\) we define the $\gamma$-adapted positive forward contribution
\begin{align}
    \tilde{z}^{(l+1)}_{ij} \coloneqq a^{(l)}_j \left( W^{(l+1)}_{ij} + \gamma \max\left\{W^{(l+1)}_{ij},0\right\} \right),
\end{align}
with \(\gamma \ge 0\) accentuating positive contributions.
Given relevances \(R^{(l+1)}_i\) at the outputs, the \(\gamma\)-rule redistributes them to inputs \(j\) as
\begin{equation}
\label{eq:lrp-gamma}
R^{(l)}_j = \sum_i \frac{\tilde{z}^{(l+1)}_{ij}}{\sum_{j'} \tilde{z}^{(l+1)}_{ij'} + \varepsilon\,\mathrm{sign}\!\left(\sum_{j'} \tilde{z}^{(l+1)}_{ij'}\right)}\, R^{(l+1)}_i,
\end{equation}
with a small \(\varepsilon>0\) for numerical stability.

\paragraph{SplitLRP implementation details}
In the DC Decomposition, each layer consists of a positive and a negative split-stream, which we treat as two sets of neurons within a unified model. Accordingly, LRP on the split-pair produces two relevance tensors in each layer, $R^{(l,+)}$ and $R^{(l,-)}$.
Since the LRP rule ignores any activation functions in the backward pass no special care is needed for the Maxout activation functions of the split-pair.
We propagate relevance from the class-specific positive output $a^{(L,+)}$. Since all weights in the decomposed model are non-negative, we use the standard LRP-$\gamma$ rule with $\gamma=0$ and $\varepsilon = 10^{-6}$, essentially forming the so-called $\epsilon$-LRP rule. As common practice, layer-normalization modules are merged into adjacent linear or convolutional layers, and relevance is split evenly across residual addition branches.
We scale and correct the activations in the forward pass just as we do in SplitCAM.
Because LRP normalizes relevance at each layer, see \cref{eq:lrp-gamma}, the total relevance in $R^{(l,+)}$ and $R^{(l,-)}$ always sums to $a^{(L,+)}$, so no additional stabilization is required during the backward pass.

We additionally report \emph{combined} LRP maps, defined as
\begin{align}
    R^{(l,\text{comb})} \coloneqq R^{(l,+)} - R^{(l,-)}.
\end{align}
This quantity aggregates the positive and negative streams into a single saliency map motivated by the following intuition: 
$R^{(l,+)}$ highlights evidence supporting the prediction, while $R^{(l,-)}$ captures evidence against it.

\paragraph{Maxpooling Nonlinearities}
We propagate relevance backwards through Maxpooling layers using the ``winner-takes-all'' scheme, assigning all relevance to the location that achieved the maximum activation.